\section{Einleitung}
Im Rahmen des 1. Semesters wird ein Projekt durchgeführt. Die Studierenden arbeiten in Gruppen und definieren ihre Problemstellung selbst. Die verwendete Programmiersprache ist entweder Go oder Python. Individuelle Wünsche sind vorgängig abgeklärt, wie die Verwendung einer anderen Programmiersprache. 

\subsection{Ausgangslage}

Das Projekt MumbleMates ist nicht eine Lösung eines Problems, sondern eher ein Learning von Go, P2P und Terminal UI. MumbleMates ist ein P2P Chat, welches in Go entwickelt wird.\\ Ziel des Projekts ist, die Grundlagen von Go zu erlernen und eine erste Version eines P2P Chats zu erstellen. Das Projekt wird in LaTeX dokumentiert.
\newline

Beide Autoren haben einen Abschluss als Informatiker Applikationsentwicklung EFZ und mit Go durfte bereits während dem Unterricht einige Erfahrungen mit Tutorials gesammelt werden. Mit LaTeX haben beide Autoren keine Erfahrungen.

\begin{itemize}
    \item Tim Ineichen hat Erfahrungen mit React, PHP und C\# (ASP.Net Core) und durfte so viel Erfahrung mit Web-Development sammeln.
    \item Tuan Binh Tran bringt Erfahrungen in Angular, C\# (ASP.Net Core) und ebenfalls viel Erfahrung im Web-Development.
\end{itemize}




\subsection{Beruflicher Lebenslauf}

Ein kurzer Überblick über die berufliche Laufbahn des Autoren:
\newline

\textbf{Tim Ineichen}
\begin{itemize}
    \item Seit 2023: Software Developer / Project Lead Synprovis GmbH
    \item 2023-2023: Mortar Operator Swiss Armed Forces
    \item 2022-2022: Requirements Engineer Klett und Balmer Verlag
    \item 2021-2022: Full Stack Developer bei der H2G Internetagentur
    \item 2019-2021: Lehrling als Informatiker Applikationsentwicklung EFZ bei der Firma Pixel Ahoi! GmbH
    \item 2017-2021: Lehre als Informatiker Applikationsentwicklung EFZ bei der Benedict Schule
\end{itemize}

\textbf{Tuan Binh Tran}
\begin{itemize}
    \item Seit 2021: Informatiker Applikationsentwicklung EFZ bei der Firma Dialog Verwaltungs-Data AG
    \item 2017-2021: Lehre als Informatiker Applikationsentwicklung EFZ bei der Firma Dialog Verwaltungs-Data AG
\end{itemize}

\subsection{Ziele}

Folgende Ziele sind für das Projekt definiert:

\subsubsection{Go}
\begin{itemize}
    \item Wir beherrschen die Grundlagen von Go und könne auch komplexere Konzepte wie Go-Routines verstehen und anwenden.
    \item Wir erstellen ein P2P Chat in Go, welches sich auf das lokale Netzwerk beschränkt.
    \item Das Ergebnis ist eine ausführbare Anwendung für Windows.
    \item Eine funktionierende Erstversion soll innerhalb der ersten zwei Wochen entstehen.
\end{itemize}

\subsubsection{LaTeX}

\begin{itemize}
    \item Die gesamte technische Dokumentation wird in LaTeX schreiben und mit Git versioniert.
    \item Das wird schlussendlich als PDF-Datei bereitgestellt.
    \item Für die Kollaboration soll die Einrichtung von LaTeX möglichst einfach gehalten werden.
    \item Wenn LaTeX gut funktioniert, wird es auch für zukünftige Projekte verwendet. 
    \item Eine erste Version soll innerhalb der ersten Woche stehen und auf der Plattform GitHub bereitgestellt werden.
\end{itemize}

\subsection{Zeitplanung}
\begin{itemize}
    \item Individueller Zeitplan mit bilateraler Absprache
    \item Projekt muss bis zum 16.03.2025 abgeschlossen sein
\end{itemize}

\newpage
